\documentclass[11pt,a4paper]{article}
\usepackage[margin=1in]{geometry}
\usepackage{graphicx}
\usepackage{hyperref}
\usepackage{enumitem}
\usepackage{titlesec}
\usepackage{fancyhdr}

% Header and footer
\pagestyle{fancy}
\fancyhf{}
\rhead{Week 1 Report}
\lhead{Online Trackers for Small Object Detection}
\rfoot{Page \thepage}

\title{\textbf{Weekly Progress Report - Week 1} \\
	\large Improve performance of existing online trackers for small object detection [UAV Videos]}

\author{Group: The Convolutioners\\\{Aagam Sheth, Mahima Parekh, Khevn Thanki, Riddhi Bhargava\}}

\date{Reporting Period: February 16 - February 21, 2026 \\
	Submission Date: February 21, 2026}


\begin{document}
	
	\maketitle
	
	\section*{Outline of Performed Tasks:}
	
	\begin{itemize}[leftmargin=*]
		\item Reviewed reference materials provided by TA including CVPR 2025 paper, MOT metrics documentation, and GitHub repositories for Kalman Filter implementation, and BoxMOT framework
		\item Successfully downloaded and explored the AU Drone Dataset
		\item Studied challenges specific to small object detection in UAV videos such as low pixel resolution, scale variation, background clutter, motion blur, and occlusions
	\end{itemize}
	
	\section*{Project Overview}
	
	\textbf{Problem Statement:} Improve existing online trackers for small object detection in UAV videos. This study evaluates Computer Vision methods for small object tracking in UAV videos.
	
	Small object detection in UAV imagery presents unique challenges because objects often occupy a very small fraction of the frame (sometimes less than 1\% of total pixels). Due to downsampling in convolutional neural networks, fine-grained spatial features may be lost, reducing localization accuracy. Additionally, cluttered backgrounds and dense multi-object scenes increase false positives and identity switches during tracking. 
	
	The objective is to design a robust pipeline that enhances detection quality for small objects and improves tracking consistency using online trackers evaluated through standard MOT metrics.
	
	\section*{Reference Materials Studied}
	
	\begin{enumerate}[leftmargin=*]
		\item Primary research paper from CVPR 2025: \\
		\url{https://cvpr.thecvf.com/virtual/2025/poster/35174}
		
		Studied advancements in small object detection and multi-object tracking, focusing on feature representation, motion modeling, and robustness in aerial scenarios.
		
		\item MOT tracking metrics documentation: \\
		\url{https://miguel-mendez-ai.com/2024/08/25/mot-tracking-metrics}
		
		Reviewed evaluation metrics including:
		\begin{itemize}
			\item MOTA (Multiple Object Tracking Accuracy)
			\item MOTP (Multiple Object Tracking Precision)
			\item IDF1 (Identity F1 Score)
			\item False Positives (FP) and False Negatives (FN)
		\end{itemize}
		
		\item GitHub repositories for implementation reference:
		\begin{itemize}
			\item Multiple Object Tracking using Kalman Filter: \\
			\url{https://github.com/NickNair/Multiple-Object-Tracking-using-Kalman-Filter}
			
			Understood motion prediction using state-space models and covariance updates.
			
			\item BoxMOT tracker framework: \\
			\url{https://github.com/mikel-brostrom/boxmot}
			
			Explored integration of detection models (e.g., YOLO) with tracking algorithms such as DeepSORT and ByteTrack.
		\end{itemize}
	\end{enumerate}
	
	\section*{Outcomes:}
	
	\begin{itemize}[leftmargin=*]
		\item Understood classical tracking methods (SORT, Mean-Shift, CamShift) and Kalman Filtering
		\item Understood AU Drone Dataset structure including annotation format and video sequences
		\item Identified key challenges in small object detection affecting tracker performance
		\item Understood relationship between detection accuracy and tracking metrics
	\end{itemize}
	
	\section*{Tentative List of Tasks for Next Week:}
	
	\begin{enumerate}[leftmargin=*]
		\item Set up development environment (Python, OpenCV, PyTorch/TensorFlow, YOLO v8) and complete study of reference papers
		\item Understand MOT evaluation metrics (MOTA, MOTP, IDF1)
	\end{enumerate}
	
\end{document}